\chapter{The Lattice Annealing Law and the CMB Timeline}
\label{app:lattice_annealing}
\addtocontents{toc}{\protect\setcounter{tocdepth}{0}}

\noindent\textbf{Abstract:} Within standard cosmology, recombination and photon decoupling occur at a characteristic epoch of approximately $380{,}000$ years, determined from the coupled thermal, atomic, and expansion history. Tension Field Theory (TFT) instead constructs a corresponding Lattice Annealing Time ($t_{\mathrm{a}}$): the physical duration required for the high-energy lattice excitation generated by Catastrophic Decompression to cool to the adopted $T_{\mathrm{a}} \sim 3000$ K stability threshold for neutral hydrogen and optical transparency. Under the bulk-cooling similarity and geometric annealing-mode closures developed below, TFT obtains $t_{\mathrm{a}} \approx 3.86 \times 10^5$ years.


\section{The Diffusive Cooling Model}

Following the State II $\to$ State I phase transition, the immense kinetic and structural heat of decompression spreads through the nascent Face-Centered Cubic (FCC) lattice. To characterize the macroscopic bulk cooling, TFT adopts the similarity scaling
\begin{equation}
    T(t) \approx \frac{T_{\mathrm S}}{\sqrt{t/t_{\mathrm S}}}.
\end{equation}

Here $T_{\mathrm S}$ is the kinetic-equivalent temperature scale associated with the fundamental substrate phase transition rather than an ordinary photon blackbody temperature.

This $t^{-1/2}$ dependence is employed here as a bulk annealing closure for the early thermal relaxation of the active lattice. A complete microscopic FCC thermal-transport derivation would require the corresponding thermal diffusivity, characteristic length scale, initial state, and boundary conditions.

\vspace{0.5em}
\noindent\textit{Note: The $t^{-1/2}$ scaling has the same mathematical time dependence as the familiar radiation-era temperature scaling used in standard cosmology. In TFT, however, it is interpreted as a structural cooling similarity law of the static active lattice rather than as a consequence of metric expansion.}


\section{The Bulk Calculation}

We calculate the bulk time ($t_{\mathrm{bulk}}$) required for the macroscopic lattice to cool from the genesis source temperature to the adopted annealing and hydrogen-stability threshold:

\begin{itemize}
    \item \textbf{Substrate Phase Transition Temperature ($T_{\mathrm S}$):} The kinetic equivalent of the fundamental substrate tension released during decompression,
    \begin{equation}
        T_{\mathrm S} \approx 1.4168 \times 10^{32}\ \mathrm{K}.
    \end{equation}

    \item \textbf{Annealing and Transparency Temperature ($T_{\mathrm a}$):} The adopted thermodynamic threshold below which stable neutral-hydrogen formation becomes possible and the coupled radiation--matter environment approaches optical transparency,
    \begin{equation}
        T_{\mathrm a} \approx 3.0 \times 10^3\ \mathrm{K}.
    \end{equation}
    A first-principles TFT derivation of this threshold from the microscopic hydrogen-knot sector remains a separate closure requirement.

    \item \textbf{Substrate Kinematic Tick ($t_{\mathrm S}$):} The stiffness-time equivalence of the lattice,
    \begin{equation}
        t_{\mathrm S} = \sqrt{\frac{\hbar G}{c_{\mathrm E}^5}} \approx 5.3912 \times 10^{-44}\ \mathrm{s}.
    \end{equation}
    It represents the fundamental kinematic action timescale of the TFT substrate.
\end{itemize}

Substituting the target annealing temperature into the bulk-cooling relation:
\begin{equation}
    T_{\mathrm a} = \frac{T_{\mathrm S}}{\sqrt{t_{\mathrm{bulk}}/t_{\mathrm S}}}.
\end{equation}

Rearranging to solve for $t_{\mathrm{bulk}}$:
\begin{equation}
    3.0 \times 10^3 = \frac{1.4168 \times 10^{32}}{\sqrt{t_{\mathrm{bulk}} / 5.3912 \times 10^{-44}}},
\end{equation}
\begin{equation}
    (3.0 \times 10^3)^2 = \frac{(1.4168 \times 10^{32})^2}{t_{\mathrm{bulk}} / 5.3912 \times 10^{-44}}.
\end{equation}

Solving yields the bulk cooling time:
\begin{equation}
    t_{\mathrm{bulk}} \approx 1.2024 \times 10^{14}\ \mathrm{s}.
\end{equation}

Converting to years ($1\ \mathrm{yr} \approx 3.1558 \times 10^7\ \mathrm{s}$):
\begin{equation}
    t_{\mathrm{bulk}} \approx 3.8102 \times 10^6\ \mathrm{years}.
\end{equation}


\section{The Geometric Annealing-Mode Closure}

The preceding bulk calculation represents the unrestricted similarity timescale of the cooling law. The macroscopic FCC lattice, however, is a structured three-dimensional medium whose finite-domain thermal relaxation must ultimately be determined by the eigenmodes of the corresponding thermal boundary-value problem.

A complete derivation would require solving
\begin{equation}
    \frac{\partial T}{\partial t} = \chi_{\mathrm{FCC}}\nabla^2 T
\end{equation}
for the appropriate FCC thermal diffusivity $\chi_{\mathrm{FCC}}$, macroscopic relaxation geometry, and thermal boundary conditions. These microscopic quantities are not yet independently closed in the present TFT construction.

For the present macroscopic annealing correspondence, TFT therefore adopts the leading dimensionless geometric mode factor $\pi^2$ as a provisional annealing-mode closure. The resulting annealing timescale is
\begin{equation}
    t_{\mathrm a} = \frac{t_{\mathrm{bulk}}}{\pi^2},
\end{equation}
which gives
\begin{equation}
    t_{\mathrm a} \approx \frac{3.8102 \times 10^6}{9.8696} \approx 3.86 \times 10^5\ \mathrm{years}.
\end{equation}

The factor $\pi^2$ is therefore not asserted here to be a universal three-dimensional FCC Laplacian eigenvalue. It represents the present geometric annealing-mode closure and must ultimately be replaced or recovered by an explicit FCC thermal boundary-value solution.

The annealing endpoint is consequently
\begin{equation}
    T(t_{\mathrm a}) = T_{\mathrm a} \sim 3000\ \mathrm{K},
\end{equation}
so that the same pair $(T_{\mathrm a},t_{\mathrm a})$ defines the annealing/transparency boundary and the initial thermal boundary used for the subsequent Fossil-CMB evolution.

\vspace{0.5em}
\noindent\textbf{Conclusion:} Under the adopted bulk-cooling similarity law and geometric annealing-mode closure, TFT obtains a lattice annealing timescale of
\begin{equation}
    t_{\mathrm a} \approx 3.86 \times 10^5\ \mathrm{years}.
\end{equation}

This lies within approximately $2\%$ of the conventional $\sim 380{,}000$-year recombination epoch. In TFT, this correspondence is interpreted strictly as a thermodynamic cooling timescale rather than as the time required to establish the macroscopic BAO length scale. The macroscopic structural scales are associated with the near-instantaneous $c_{\mathrm S}$ longitudinal shockwaves generated during Catastrophic Decompression. The numerical agreement therefore establishes a robust, non-expanding lattice-annealing correspondence to the observed recombination epoch, while a complete first-principles prediction requires microscopic derivations of both the hydrogen stability threshold and the FCC thermal transport eigenmode.
